\section{Marco teórico}
\label{sec:marco}

\subsection{Análisis de sentimiento}

El análisis de sentimiento, también llamado minería de opiniones, es el campo que estudia
la extracción automática de actitudes, opiniones y emociones a partir de texto
\citep{liu2012sentiment}. La literatura distingue dos grandes familias de métodos: los
enfoques supervisados, que entrenan clasificadores sobre corpus etiquetados, y los
enfoques léxicos, que combinan un diccionario de términos con una regla de agregación
\citep{taboada2011socal}. Los métodos léxicos, pese a su simplicidad, presentan ventajas
notables en contextos de dominio específico y de escasez de datos etiquetados: son
transparentes, interpretables y no requieren entrenamiento.

Ejemplos representativos de la familia léxica son SO-CAL \citep{taboada2011socal} y
VADER \citep{hutto2014vader}. Este último fue dise\~nado específicamente para texto de
redes sociales e incorpora reglas para el manejo de emojis, mayúsculas, puntuación y
grado de amplificación. El método implementado en este trabajo sigue esa tradición, pero
se adapta al espa\~nol y se extiende con un esquema multietiqueta de emociones.

\subsection{Enfoques léxicos en espa\~nol}

El espa\~nol ha recibido atención creciente en la literatura de sentimiento.
\citet{perez2012spanish} construyeron léxicos de sentimiento en espa\~nol mediante
traducción y validación de recursos en inglés, mientras que \citet{mohammad2013emolex}
desarrollaron el léxico NRC de asociación palabra--emoción, posteriormente adaptado a
múltiples idiomas. Estos recursos subrayan un problema central: la traducción directa
introduce ruido y pierde matices culturales. Por ello, un diccionario específico del
dominio, construido y revisado manualmente, suele superar a los recursos genéricos en
tareas aplicadas.

El diccionario de este proyecto consta de 320 términos ponderados, organizados en 18
categorías emocionales y una escala de polaridad que va de $-3$ a $+5$. Cada término
aporta simultáneamente una contribución a la polaridad y una o varias contribuciones
emocionales, lo que permite el etiquetado multietiqueta.

\subsection{Tratamiento de la negación}

La negación es uno de los fenómenos que más degrada el rendimiento de los métodos
léxicos. Una frase como ``\emph{no es buen\'\i simo}'' contiene un término positivo
(``buen\'\i simo'') cuyo sentido se invierte por el negador. La literatura ha abordado
el problema mediante reglas de alcance \citep{brooke2009negation} o modelos de
meta-aprendizaje \citep{morante2012negation}. La solución adoptada aquí es una ventana
contextual de tres tokens: si un término léxico aparece precedido por un negador dentro
de esa ventana, su polaridad se invierte. Los negadores reconocidos incluyen ``no'',
``ni'', ``nunca'', ``jam\'as'', ``tampoco'', ``sin'', ``nada'', ``nadie'' y variantes
coloquiales como ``jama''.

\subsection{Emojis y reacciones como señal}

Los emojis no son ruido: constituyen una forma paralela de comunicación afectiva. El
estudio de referencia \citet{novak2015emojis} asignó puntuaciones de sentimiento a más
de mil emojis y demostró que su distribución es sistemática y culturalmente consistente.
En consecuencia, este trabajo \emph{nunca} elimina los emojis: los procesa como una
señal independiente que contribuye a la polaridad y a las emociones. Análogamente, las
reacciones (los \emph{stickers} con contador) se registran como una tercera señal,
aunque por defecto se les asigna peso cero para evitar doble conteo, dado que una
reacción no siempre se refiere al mensaje inmediato.

\subsection{Modelos de emociones}

La psicología ha propuesto varios modelos de la estructura emocional. El modelo de
emociones básicas de \citet{ekman1992emotions} postula un conjunto reducido de
emociones universales (alegría, tristeza, miedo, ira, sorpresa y asco). El modelo
psicoevolucionario de \citet{plutchik1980emotions} organiza las emociones en una rueda
de ocho dimensiones con intensidades graduadas. El modelo circunflejo de
\citet{russell1980circumplex} representa el afecto en dos ejes ortogonales: valencia
(agrado--desagrado) y activación (activación--desactivación). El esquema multietiqueta
de dieciocho categorías empleado aquí es ecléctico: recoge emociones básicas, matices
sociales (gratitud, confianza, empatía) y estados relacionados con la tarea académica
(ansiedad ante plazos, frustración procedimental), inspirados en la rueda de Plutchik y
en el léxico NRC.

\subsection{Modelado de tópicos}

Para descubrir la estructura temática sin supervisión se utiliza la asignación latente
de Dirichlet (LDA), introducida por \citet{blei2003lda} y popularizada por
\citet{blei2012topicmodels}. LDA representa cada documento como una mezcla de tópicos y
cada tópico como una distribución sobre el vocabulario. La advertencia metodológica de
\citet{grimmer2013text} es pertinente: los tópicos son construcciones estadísticas que
deben interpretarse con cautela y validarse cualitativamente, no como categorías
objetivas preexistentes.

\subsection{Cadenas de Markov}

Para modelar la dinámica de los estados emocionales se emplean cadenas de Markov, un
formalismo clásico para procesos estocásticos sobre un espacio de estados finito
\citep{norris1997markov}. Cada mensaje se asigna a un estado (positivo, negativo o
neutro) y se estima la matriz de transición a partir de las frecuencias observadas. La
distribución estacionaria describe la proporción de tiempo que el sistema pasa en cada
estado a largo plazo, y el tiempo medio de primer retorno o de permanencia caracteriza
la estabilidad de cada estado. Este análisis añade una dimensión dinámica que la mera
distribución estática de polaridad no captura.
